\documentclass[12pt]{article}
\usepackage[utf8]{inputenc}
\usepackage[T1]{fontenc}
\usepackage{hyperref}
\title{Wamin Language Resources for K-12 Curriculum Development}
\author{Compiled by Claude Sonnet 5.5 (claude-sonnet-5-5) for Andr\'as Kornai}
\date{}
\begin{document}
\maketitle

\section{Introduction}
Wamin (ISO 639-3 wmi; Glottocode wami1239), also called Agwamin or Ewamian, is a Pama-Nyungan (Southern Paman) language of North Queensland. It was historically centred on the Etheridge region, around Einasleigh, Georgetown and Mount Surprise. Wikipedia lists it as extinct by 2005, with a revival reported underway. One source mentions a single living speaker in 1981. Search results showed very little documentation of its lexicon or grammar. The AIATSIS record (Y132) indicates little or none.

\section{Primary Educational Resources}
None were found in search results.

\section{Linguistic and Reference Materials}
\begin{itemize}
\item Wikipedia, Wamin language: \url{https://en.wikipedia.org/wiki/Wamin_language}
\item AIATSIS Austlang record Y132 (Agwamin): \url{https://collection.aiatsis.gov.au/austlang/language/Y132?code=Y132}
\item Wikipedia, Wakaman, which is sometimes treated as a variant of or overlapping with Agwamin: \url{https://en.wikipedia.org/wiki/Wakaman}
\end{itemize}
Peter Sutton's early 1970s fieldwork with elder Fred Fulford is mentioned in a source. I saw no URL for the primary materials.

\section{Religious and Community Resources}
None were found in search results. A revival effort is mentioned, but I found no organisation or URL.

\section{Digital and Audio Resources}
None were found in search results.

\section{Curriculum Development Context}
There are no known school programmes. A heritage and revitalization curriculum would depend on archival materials. A sensible next step is to approach AIATSIS and the Ewamian people directly. Because this is a revitalization context, community ownership of any materials is essential. The Agwamin/Wamin/Wakaman relationship should be clarified with the community.

\end{document}